\section{Setup}
\label{section-setup}

\noindent
We will throughout work over a locally Noetherian universally
catenary base $S$ endowed with a dimension function $\delta$.
Although it is likely possible to generalize (parts of) the
discussion in the chapter, it seems that this is a good first
approximation. It is exactly the generality discussed in \cite{Thorup}.
We usually do not assume our schemes are
separated or quasi-compact. Many interesting algebraic stacks
are non-separated and/or non-quasi-compact and this is a good
case study to see how to develop a reasonable theory for those as well.
In order to reference these hypotheses we give it a number.

\begin{situation}
\label{situation-setup}
Here $S$ is a locally Noetherian, and universally catenary scheme.
Moreover, we assume $S$ is endowed with a dimension function
$\delta : S \longrightarrow \mathbf{Z}$.
\end{situation}

\noindent
See Morphisms, Definition \ref{morphisms-definition-universally-catenary}
for the notion of a universally catenary scheme, and see
Topology, Definition \ref{topology-definition-dimension-function}
for the notion of a dimension function. Recall that any locally
Noetherian catenary scheme locally has a dimension function, see
Properties, Lemma \ref{properties-lemma-catenary-dimension-function}.
Moreover, there are lots of schemes which are universally catenary,
see Morphisms, Lemma \ref{morphisms-lemma-ubiquity-uc}.

\medskip\noindent
Let $(S, \delta)$ be as in Situation \ref{situation-setup}.
Any scheme $X$ locally of finite type over $S$ is locally Noetherian
and catenary. In fact, $X$ has a canonical dimension function
$$
\delta = \delta_{X/S} : X \longrightarrow \mathbf{Z}
$$
associated to $(f : X \to S, \delta)$ given by the rule
$\delta_{X/S}(x) = \delta(f(x)) + \text{trdeg}_{\kappa(f(x))}\kappa(x)$.
See Morphisms, Lemma \ref{morphisms-lemma-dimension-function-propagates}.
Moreover, if $h : X \to Y$ is a morphism of schemes locally of finite
type over $S$, and $x \in X$, $y = h(x)$,
then obviously
$\delta_{X/S}(x) = \delta_{Y/S}(y) + \text{trdeg}_{\kappa(y)}\kappa(x)$.
We will freely use this function and its properties in the following.

\medskip\noindent
Here are the basic examples of setups as above.
In fact, the main interest lies in the case where the base
is the spectrum of a field, or the case where the base
is the spectrum of a Dedekind ring (e.g.\ $\mathbf{Z}$,
or a discrete valuation ring).

\begin{example}
\label{example-field}
Here $S = \Spec(k)$ and $k$ is a field.
We set $\delta(pt) = 0$ where $pt$ indicates the unique point of $S$.
The pair $(S, \delta)$ is an example of a situation as in
Situation \ref{situation-setup} by
Morphisms, Lemma \ref{morphisms-lemma-ubiquity-uc}.
\end{example}

\begin{example}
\label{example-domain-dimension-1}
Here $S = \Spec(A)$, where $A$ is a Noetherian domain
of dimension $1$.
For example we could consider $A = \mathbf{Z}$.
We set $\delta(\mathfrak p) = 0$ if
$\mathfrak p$ is a maximal ideal and $\delta(\mathfrak p) = 1$
if $\mathfrak p = (0)$ corresponds to the generic point.
This is an example of Situation \ref{situation-setup} by
Morphisms, Lemma \ref{morphisms-lemma-ubiquity-uc}.
\end{example}

\begin{example}
\label{example-CM-irreducible}
Here $S$ is a Cohen-Macaulay scheme. Then $S$ is universally catenary by
Morphisms, Lemma \ref{morphisms-lemma-ubiquity-uc}.
We set $\delta(s) = -\dim(\mathcal{O}_{S, s})$.
If $s' \leadsto s$ is a nontrivial specialization of points of $S$,
then $\mathcal{O}_{S, s'}$ is the localization of $\mathcal{O}_{S, s}$
at a nonmaximal prime ideal $\mathfrak p \subset \mathcal{O}_{S, s}$, see
Schemes, Lemma \ref{schemes-lemma-specialize-points}.
Thus $\dim(\mathcal{O}_{S, s}) =
\dim(\mathcal{O}_{S, s'}) + \dim(\mathcal{O}_{S, s}/\mathfrak p) >
\dim(\mathcal{O}_{S, s'})$ by
Algebra, Lemma \ref{algebra-lemma-CM-dim-formula}.
Hence $\delta(s') > \delta(s)$. If $s' \leadsto s$ is
an immediate specialization, then there is no prime
ideal strictly between $\mathfrak p$ and $\mathfrak m_s$
and we find $\delta(s') = \delta(s) + 1$. Thus $\delta$
is a dimension function. In other words, the pair $(S, \delta)$
is an example of Situation \ref{situation-setup}.
\end{example}

\noindent
If $S$ is Jacobson and $\delta$ sends closed points to zero, then $\delta$
is the function sending a point to the dimension of its closure.

\begin{lemma}
\label{lemma-delta-is-dimension}
Let $(S, \delta)$ be as in Situation \ref{situation-setup}.
Assume in addition $S$ is a Jacobson scheme, and $\delta(s) = 0$ for every
closed point $s$ of $S$. Let $X$ be locally of finite type over $S$.
Let $Z \subset X$ be an integral closed subscheme and let
$\xi \in Z$ be its generic point. The following integers are the same:
\begin{enumerate}
\item $\delta_{X/S}(\xi)$,
\item $\dim(Z)$, and
\item $\dim(\mathcal{O}_{Z, z})$ where $z$ is a closed point of $Z$.
\end{enumerate}
\end{lemma}

\begin{proof}
Let $X \to S$, $\xi \in Z \subset X$ be as in the lemma.
Since $X$ is locally of finite type over $S$ we see that
$X$ is Jacobson, see
Morphisms, Lemma \ref{morphisms-lemma-Jacobson-universally-Jacobson}.
Hence closed points of $X$ are dense in every closed subset of $Z$
and map to closed points of $S$. Hence given any chain
of irreducible closed subsets of $Z$ we can end it with a closed point of $Z$.
It follows that $\dim(Z) = \sup_z(\dim(\mathcal{O}_{Z, z})$
(see Properties, Lemma \ref{properties-lemma-codimension-local-ring})
where $z \in Z$ runs over the closed points of $Z$.
Note that $\dim(\mathcal{O}_{Z, z}) = \delta(\xi) - \delta(z)$
by the properties of a dimension function.
For each closed $z \in Z$ the field extension
$\kappa(z)/\kappa(f(z))$ is finite, see Morphisms,
Lemma \ref{morphisms-lemma-jacobson-finite-type-points}.
Hence $\delta_{X/S}(z) = \delta(f(z)) = 0$ for $z \in Z$ closed.
It follows that all three integers are equal.
\end{proof}

\noindent
In the situation of the lemma above the
value of $\delta$ at the generic point of a closed irreducible subset
is the dimension of the irreducible closed subset.
However, in general we cannot expect the equality to hold.
For example if $S = \Spec(\mathbf{C}[[t]])$ and
$X = \Spec(\mathbf{C}((t)))$ then we would get
$\delta(x) = 1$ for the unique point of $X$, but $\dim(X) = 0$.
Still we want to think of $\delta_{X/S}$ as giving the
dimension of the irreducible closed subschemes. Thus we introduce
the following terminology.

\begin{definition}
\label{definition-delta-dimension}
Let $(S, \delta)$ as in Situation \ref{situation-setup}.
For any scheme $X$ locally of finite type over $S$
and any irreducible closed subset $Z \subset X$ we define
$$
\dim_\delta(Z) = \delta(\xi)
$$
where $\xi \in Z$ is the generic point of $Z$.
We will call this the {\it $\delta$-dimension of $Z$}.
If $Z$ is a closed subscheme of $X$, then we define
$\dim_\delta(Z)$ as the supremum of the $\delta$-dimensions
of its irreducible components.
\end{definition}







